\section{Use the included LaTeX template}
\label{sec:template}

The source for this guide retains the preamble used by
OGPO~\citep{patil2026ogpo}. It is a working paper archive, not a generic LaTeX
package. That distinction matters. The template contains reusable theorem,
reference, color, and editing machinery, but it also contains OGPO-specific
algorithm names and notation. Keep the former. Replace the latter before
starting a new paper.

Begin with a clean copy of the project. Replace the files in \path{body/} and
\path{appendix/}, clear the bibliography, and update the title block in
\path{latex-template-main.tex}. Then remove unused definitions from
\path{specific_macros.tex} and \path{algnames.tex}. A copied macro that never
appears in the new paper is not harmless: it makes collisions and stale
terminology harder to find.

\subsection{Assign each choice to one file}

The include order lives in \path{preamble/_preamble_includes.tex}. Keep that
order until there is a specific reason to change it. Each file controls a
different part of the build:

\begin{description}[leftmargin=*,style=nextline]
    \item[\texttt{color\_defs.tex}]
    Defines link colors, theorem colors, emphasis colors, and the
    \texttt{AIbox} environment. It loads before files that use those colors.

    \item[\texttt{header.tex}]
    Loads common packages and defines page geometry, citations, cross
    references, algorithms, and the arXiv-dependent heading commands.

    \item[\texttt{theorem\_styles.tex}]
    Defines theorem-like environments and their appearance.

    \item[\texttt{characters.tex} and \texttt{general\_macros.tex}]
    Define reusable mathematical characters, delimiters, operators, and
    formatting commands.

    \item[\texttt{specific\_macros.tex}]
    Holds notation that belongs to one paper. This file should change
    substantially when the project changes.

    \item[\texttt{algnames.tex}]
    Defines method names, method colors, task names, and the algorithm font.
    Treat it as a paper-specific legend.

    \item[\texttt{color-edits.sty} and \texttt{comment\_macros.tex}]
    Define temporary author edits, comments, deletion marks, and drafting
    flags.

    \item[\texttt{arxiv\_title.tex}]
    Defines the public title layout and metadata commands for websites, code,
    documentation, models, and dates.
\end{description}

The preamble is easiest to maintain when each new command goes into the
narrowest file whose purpose matches the command. A task name does not belong
in \texttt{header.tex}. A package import does not belong in
\texttt{specific\_macros.tex}.

\subsection{Use the theorem environments that are actually defined}

With \texttt{customthms} enabled, choose an environment by the kind of
statement:

\begin{itemize}[leftmargin=*]
    \item Use \texttt{theorem}, \texttt{lemma}, \texttt{corollary},
    \texttt{proposition}, or \texttt{claim} for results.
    \item Use \texttt{definition}, \texttt{assumption}, or
    \texttt{condition} for definitions and conditions.
    \item Use \texttt{remark} or \texttt{example} for commentary and
    examples.
\end{itemize}

The template defines several less common environments as well. Use the full
names \texttt{definition} and \texttt{theorem}; it does not define
\texttt{defn} or \texttt{thm}.

\begin{verbatim}
\begin{definition}[Stable rollout]\label{def:stable-rollout}
A rollout is stable if ...
\end{definition}

\begin{theorem}[Main guarantee]\label{thm:main}
Under \Cref{def:stable-rollout}, ...
\end{theorem}

By \Cref{thm:main}, the error ...
\end{verbatim}

Put the label immediately after the opening environment. Use
\texttt{\textbackslash Cref} at the beginning of a sentence and
\texttt{\textbackslash cref} within one. The template already names equations,
assumptions, conditions, and components for \texttt{cleveref}. Equations are
numbered by section.

Starred forms such as \texttt{theorem*}, \texttt{definition*}, and
\texttt{remark*} suppress numbering. Use them only when no later sentence
will need to cite the statement. \texttt{thminformal} provides a numbered
informal theorem. An informal statement should still match the assumptions
and conclusion of its formal counterpart.

The template also contains helpers for lettered variants of an existing
result. For example:

\begin{verbatim}
\begin{thmmod}{thm:main}{b}{Finite-sample form}
...
\end{thmmod}
\end{verbatim}

This prints a theorem number based on \texttt{thm:main} with the suffix
\texttt{b}. Related helpers include \texttt{lemmod}, \texttt{propmod},
\texttt{defnmod}, \texttt{asmmod}, \texttt{cormod}, and \texttt{condmod}.
Use them only when the new statement is genuinely a named variant of an
existing result. Several other inherited helpers assume definitions from the
original project and are not a stable public interface. Compile a minimal
example before relying on one.

When \texttt{customthms} is false, the template does not define the full set of
environments. That branch assumes the venue style provides them. Do not turn
the toggle off in a plain \texttt{article} build and expect an identical
theorem system.

\subsection{Choose build toggles by their effect}

The toggles appear near the top of \texttt{latex-template-main.tex}:

\begin{verbatim}
\newtoggle{arxiv}
\toggletrue{arxiv}
\newtoggle{customthms}
\toggletrue{customthms}
\end{verbatim}

\texttt{arxiv} controls the public title, geometry, fonts, page style,
algorithm packages, and colored paragraph heads. The false branch is a hook
for a venue build, not a complete venue style by itself. In particular, the
template loads \texttt{algorithm} and \texttt{algorithmic} only in the arXiv
branch. A conference style must supply compatible algorithm support before
the same source will compile.

\texttt{customthms} controls the colored theorem definitions. The
\texttt{toreturn} toggle controls the visibility of return notes; those
notes use \texttt{\textbackslash RETURNN}. The inherited \texttt{icml} and
\texttt{colt} toggles should remain only if the new paper still has those
build modes. Remove a dead toggle rather than leaving a branch no one compiles.

The arXiv title macros also exist only when the arXiv title file is loaded.
Calls such as \texttt{\textbackslash paperdate} and direct access to
\texttt{\textbackslash metadatalist} must therefore stay inside the public
build or receive venue-side definitions. The same rule applies to
\texttt{\textbackslash fancyhead}, because \texttt{fancyhdr} is loaded in the
arXiv branch.

\subsection{Use paragraph and emphasis commands deliberately}

Three commands make emphasis change with the build mode:

\begin{description}[leftmargin=*,style=nextline]
    \item[\texttt{\textbackslash colorpar\{Heading.\}}]
    Creates a paragraph heading in both modes. The heading is colored in the
    arXiv build and plain in the venue branch.

    \item[\texttt{\textbackslash togglepar\{Heading.\}}]
    Calls \texttt{colorpar} in the arXiv build, but becomes inline bold text in
    the venue branch. Use it for a compact run-in heading that can lose the
    paragraph command under a tight style.

    \item[\texttt{\textbackslash colorbold\{Text\}}]
    Produces colored bold text in the arXiv build and ordinary bold text in the
    venue branch.
\end{description}

\begin{verbatim}
\togglepar{Experimental setup.} We compare ...

\colorpar{Failure under distribution shift.}
The policy fails when ...

\colorbold{Main finding.} The intervention ...
\end{verbatim}

The period belongs inside the heading argument. Do not use these commands to
make every paragraph look important. A paragraph heading should expose a real
change in purpose: setup, result, mechanism, boundary, or implication.

\texttt{\textbackslash termbold} and
\texttt{\textbackslash takeawaybold} use fixed semantic colors rather than a
build toggle. They are appropriate only when the chosen venue permits that
color. The paper must remain understandable when all emphasis is printed in
grayscale.

\subsection{Define method names and colors once}

\texttt{algnames.tex} pairs a typographic method name with a stable color.
The base helpers are:

\begin{verbatim}
\newcommand{\algofont}[1]{%
  {\scalefont{1.0}{\texttt{\textbf{#1}}}}}
\newcommand{\algofonttitle}[1]{%
  {\scalefont{1.08}{\textbf{\texttt{#1}}}}}
\end{verbatim}

For a new method, define the color, a neutral name, a colored prose name, and
a title form:

\begin{verbatim}
\definecolor{methodcol}{HTML}{A60000}
\newcommand{\Methodcolor}{methodcol}
\newcommand{\Methodnc}{\algofont{METHOD}}
\newcommand{\Method}{{\color{\Methodcolor}\Methodnc}}
\newcommand{\Methodtitle}{%
  {\color{\Methodcolor}\algofonttitle{METHOD}}}
\end{verbatim}

Use \texttt{\textbackslash Method\{\}} in prose so the following space is not
swallowed. Use the neutral form when color would confuse a comparison or when
the venue requires black text. Reserve the title form for a title or compact
display; its font is deliberately larger.

Keep the same method-to-color assignment in plots and diagrams. Store the hex
value in the plotting code as well as the LaTeX source. Color should help a
reader track a method, but the method name, line style, or marker must still
identify it without color.

The inherited file contains several aliases for the same method. It does not
apply the neutral, colored, and title forms uniformly. Do not copy that
inconsistency into a new method. Pick one naming pattern and remove aliases
unused by the paper.

\subsection{Keep colors semantic}

\texttt{color\_defs.tex} defines four semantic color groups. Link and citation
colors use dark teal. Theorem headings use dark red. Term and takeaway macros
use stronger red accents. Method colors live in \texttt{algnames.tex}.

Change the semantic definition instead of scattering raw color commands
through the manuscript. For example, changing
\texttt{\textbackslash thmcolordark} should update every theorem heading. It
should not require a search through theorem prose.

The inherited names are not all standard color names. Inspect the actual
\texttt{\textbackslash definecolor} value before reusing one. A familiar name
can point to a project-specific hex code.

The environment named \texttt{AIbox} is an inherited
\texttt{tcolorbox}; its name does not prescribe its content. It is used for
compact summaries:

\begin{verbatim}
\begin{AIbox}{Core insight}
The critic supplies ...
\end{AIbox}
\end{verbatim}

Use a box only when the same short statement should remain visible while the
reader scans the paper. Do not place a second box inside it, and do not use a
box to repeat a contribution already stated in a heading and figure.

\subsection{Use color edits as temporary review state}

\texttt{comment\_macros.tex} registers each author:

\begin{verbatim}
\addauthor{ms}{magenta}
\addauthor{xy}{blue}
\end{verbatim}

Each registration creates four commands:

\begin{description}[leftmargin=*,style=nextline]
    \item[\texttt{\textbackslash xyedit\{new text\}}]
    Shows the text in the author's color. With the
    \texttt{suppress} option, the text remains and loses its color.

    \item[\texttt{\textbackslash xycomment\{question\}}]
    Shows a bracketed author comment. The \texttt{suppress} option hides it.

    \item[\texttt{\textbackslash xymargincomment\{question\}}]
    Adds a colored marker and small comment at that point. Despite the inherited
    name, the current implementation does not place the full text in a
    conventional margin note.

    \item[\texttt{\textbackslash xydelete\{old text\}}]
    Marks a deletion. The \texttt{showdeletions} option displays struck-out
    text; the default displays a marker and deletion note.
\end{description}

Set package options where \texttt{color-edits.sty} is loaded:

\begin{verbatim}
\usepackage[showdeletions]{preamble/color-edits}
% or, for a clean diagnostic build:
\usepackage[suppress]{preamble/color-edits}
\end{verbatim}

A suppressed build is a diagnostic, not a final cleanup. It can hide comments,
deletion markers, and return notes that still remain in the source. Before
release, remove resolved edit commands and search for author prefixes,
\texttt{\textbackslash REF}, \texttt{\textbackslash CITE},
\texttt{\textbackslash TODO}, and \texttt{\textbackslash RETURNN}.
\texttt{\textbackslash ignore} removes its argument in every build and should
not hold text that might need to return.

\subsection{Use the algorithm environment with its matching package}

The arXiv branch uses the older \texttt{algorithmic} command set. Its commands
are uppercase:

\begin{verbatim}
\begin{algorithm}[t]
\caption{Method update}\label{alg:update}
\begin{algorithmic}[1]
\REQUIRE Initial policy and replay buffer
\FOR{each update}
  \STATE Sample a batch
  \STATE Update the critic
\ENDFOR
\end{algorithmic}
\end{algorithm}
\end{verbatim}

\texttt{\textbackslash algcomment} provides the blue monospace comments used
in the OGPO pseudocode. Use it only for short operational annotations. A
paragraph of explanation belongs in the caption or surrounding text.

Do not mix \texttt{algorithmic}, \texttt{algpseudocode}, and
\texttt{algorithm2e} commands. They expose similar names with different
semantics. If a venue supplies its own algorithm package, adapt the source once
and compile both modes.

\subsection{Use notation shortcuts without hiding the mathematics}

\path{characters.tex} creates families of common symbols. Use
\texttt{\textbackslash bA} for a bold matrix and
\texttt{\textbackslash bx} for a bold vector. The same file provides
\texttt{\textbackslash bbR} for blackboard-bold \(R\) and
\texttt{\textbackslash cF} for calligraphic \(F\).

\path{general_macros.tex} defines paired delimiters. These include
\texttt{\textbackslash abs}, \texttt{\textbackslash prn}, and
\texttt{\textbackslash nrm}. It also defines operators such as
\texttt{\textbackslash argmin} and \texttt{\textbackslash argmax}.

These generated names are convenient but broad. Check for an existing command
before adding a new one. Do not redefine a familiar LaTeX command merely to
save one character. When a symbol has paper-specific meaning, define a
semantic name in \texttt{specific\_macros.tex} instead of relying on a visual
alias throughout the manuscript.

\subsection{Populate the public title metadata carefully}

The public title file can collect metadata below the title card:

\begin{verbatim}
\paperwebsite{https://project.example}
\papercode{https://github.com/group/project}
\paperdocs{https://project.example/docs}
\papermodels{https://huggingface.co/group/project}
\paperdate{\today}
\end{verbatim}

\texttt{papercode}, \texttt{paperblog}, and \texttt{papermodels} include logo
files from \texttt{logos/}. Keep those assets in the source bundle or replace
the commands with plain links. Do not advertise an empty repository, private
artifact, or placeholder URL.

The title card prints the entries in the order the commands are called. Include
only metadata that helps a reader verify, reproduce, or use the work. The venue
build should not depend on these public-title commands unless its style defines
them.

\subsection{Audit the template before submission}

Build the public and venue modes from a clean directory. For each mode, check:

\begin{enumerate}[leftmargin=*]
    \item every theorem environment exists and every cross-reference resolves;
    \item paragraph heads have the intended spacing after the toggle changes;
    \item algorithm floats compile under the selected package;
    \item method names retain the same spelling and color assignment;
    \item no author edit, deletion marker, return note, or placeholder remains;
    \item public metadata points to released artifacts; and
    \item the source bundle contains its logos, figures, bibliography, style
    files, and any venue-owned class or bibliography style.
\end{enumerate}

Finally, remove unused packages and source-specific macros. The preserved OGPO
preamble is useful because its decisions remain visible. It becomes a burden
when a new paper carries all of those decisions without a clear need for them.
